\chapter{Compréhension des Données}

Ce chapitre correspond à la phase de compréhension des données du cycle CRISP-DM. Elle couvre quatre tâches : la collecte des données, leur description structurelle, leur exploration statistique et la vérification de leur qualité. Le projet mobilise deux corpus distincts : des données de surveillance opérationnelle SAP extraites en temps réel et stockées dans MongoDB, et une base de connaissances documentaire construite à partir de 297 livres de formation SAP et du portail d'aide en ligne, indexée dans PostgreSQL avec pgvector.

% ─────────────────────────────────────────────────────────────────────────────
\section{Collecte des Données}
% ─────────────────────────────────────────────────────────────────────────────

\subsection{Données de Surveillance SAP — Extraction via PyRFC}

Les données de surveillance opérationnelle sont extraites automatiquement des systèmes SAP de production par la couche \textbf{PyRFC}, le connecteur Python pour les Remote Function Calls SAP. PyRFC permet d'appeler directement les modules de fonction SAP depuis Python sans passer par la SAP GUI, rendant l'extraction répétable, planifiable et traçable.

Chaque extraction produit des enregistrements horodatés portant un champ \texttt{SID} qui identifie le système SAP source et un champ \texttt{Extraction\_Timestamp} qui permet l'historisation. Dix-huit transactions SAP sont couvertes : journaux d'erreurs ABAP (ST22), statuts de jobs (SM37), espace base de données (DB02), entrées de verrouillage (SM12), processus de travail (SM50, SM66), files qRFC (SMQ1, SMQ2), messagerie (SOST), ainsi que les métriques système CPU, mémoire et système de fichiers.

\subsection{Base de Connaissances RAG — Extraction Multi-Agent}

La base de connaissances du chatbot est construite à partir de deux sous-sources documentaires.

La première est constituée de \textbf{297 supports de cours SAP au format PDF} : administration Basis (ADM100, ADM505), développement ABAP (BC100, BC400), modules fonctionnels MM, SD, FI, et formations spécialisées NetWeaver, HANA, Fiori, sécurité. Ces documents issus du programme de formation accrédité SAP constituent la référence pédagogique la plus complète pour un technicien SAP.

La seconde est le \textbf{portail d'aide en ligne SAP} (help.sap.com), dont des sections thématiques ont été extraites par web scraping. Ce portail fournit une documentation orientée configuration, codes d'erreur, paramètres système et notes de correction (SAP Notes), complémentaire aux livres de formation.

\paragraph{Extraction PDF par agents parallèles.}
Le script \texttt{Sap\_batch\_extractor.py} orchestre un ensemble d'agents d'extraction, chacun responsable d'un sous-ensemble de livres. Chaque agent opère de façon autonome : il charge le PDF avec \textbf{pdfplumber}, détecte la hiérarchie des sections à partir de la taille et de la graisse de police, extrait le texte de chaque section, isole les figures et produit un fichier \texttt{structured\_content.json}. Le traitement en parallèle réduit le temps total d'extraction de plusieurs heures à quelques dizaines de minutes.

\paragraph{Web scraping du portail SAP.}
L'extraction des pages d'aide utilise une chaîne de trois outils : \textbf{Scrapy} pour la navigation et le téléchargement des pages statiques, \textbf{Playwright} pour les composants dynamiques chargés en JavaScript, et \textbf{BeautifulSoup} pour le parsing et le nettoyage du HTML brut.

\paragraph{Résilience SSL lors de l'ingestion.}
L'ingestion des vecteurs dans la base PostgreSQL distante hébergée sur Supabase expose le pipeline à des interruptions de connexion SSL sur des sessions de longue durée. Deux mécanismes ont été implémentés : des TCP keepalives (\texttt{keepalives\_idle=30\,s, keepalives\_interval=10\,s}) pour maintenir la connexion active, et une boucle de retry automatique en cas de drop SSL.

\begin{lstlisting}[language=Python, caption={Boucle de retry SSL dans ingest\_to\_db.py}]
for attempt in range(3):
    try:
        conn = psycopg2.connect(
            DB_URL, keepalives=1,
            keepalives_idle=30, keepalives_interval=10
        )
        # ingestion du batch
        break
    except Exception as exc:
        if "SSL" in str(exc):
            print(f"[WARN] SSL drop, retry {attempt+1}/3")
            time.sleep(2)
        else:
            raise
\end{lstlisting}

% ─────────────────────────────────────────────────────────────────────────────
\section{Description des Données}
% ─────────────────────────────────────────────────────────────────────────────

\subsection{Collections MongoDB — Données de Surveillance}

Les données de surveillance sont réparties dans dix-huit collections MongoDB, chacune correspondant à une transaction SAP ou à une métrique système. Le tableau suivant présente les collections et leurs champs principaux.

\begin{longtable}{|p{2.4cm}|p{4.0cm}|p{7.2cm}|}
\caption{Collections MongoDB et champs principaux}
\label{tab:mongo-collections}\\
\hline
\textbf{Collection} & \textbf{Transaction SAP} & \textbf{Champs principaux} \\
\hline
\endfirsthead
\hline
\textbf{Collection} & \textbf{Transaction SAP} & \textbf{Champs principaux} \\
\hline
\endhead
\texttt{ST22}      & Erreurs d'exécution ABAP     & Current\_Date, Runtime\_Error\_Name, Application\_Server, User\_Name \\
\hline
\texttt{SM37}      & Jobs en arrière-plan          & Job, Status, Start\_Date, End\_Date, Job\_Created\_By \\
\hline
\texttt{DB02}      & Espace base de données        & TableSpace\_Name, Used\_Percent, Free\_MB, AutoExtend \\
\hline
\texttt{SM12}      & Entrées de verrouillage       & Table\_Name, User\_name, Date, Time \\
\hline
\texttt{SM51}      & Serveurs d'application        & HOST, STATE \\
\hline
\texttt{SM66}      & Processus de travail          & WP\_Type, Total, Running, WP\_Ratio \\
\hline
\texttt{SM58}      & RFC transactionnel             & Function\_Module, Target\_System, Status \\
\hline
\texttt{SM13}      & Enregistrements de mise à jour & User, TCode, Status \\
\hline
\texttt{SP01}      & Demandes de spool              & Spool\_No, User\_Name, Status \\
\hline
\texttt{ST13}      & Chaînes de processus           & Chain\_ID, Log\_ID, Status \\
\hline
\texttt{ST03N}     & Analyse de charge              & Task\_Type, Average\_Response\_Time\_per\_Step \\
\hline
\texttt{SMQ1/SMQ2} & Files qRFC                    & Queue\_Name, Destination, Status \\
\hline
\texttt{SOST}      & Messagerie sortante            & Creator, ErrorMsg \\
\hline
\texttt{CPU}       & Métriques processeur           & user\_cpu\_percent, system\_cpu\_percent, idle\_cpu\_percent \\
\hline
\texttt{MEMORY}    & Métriques mémoire              & MIN\_FR\_MEM, MAX\_FR\_MEM, AVG\_SWAPFR \\
\hline
\texttt{FSYS}      & Système de fichiers            & filesystem\_path, capacity\_gb, free\_gb \\
\hline
\texttt{MSS\_DB}   & Base SQL Server                & DB\_SIZE, LOG\_SIZE, RESERVED \\
\hline
\end{longtable}

Chaque enregistrement porte également les champs transversaux \texttt{SID} (identifiant du système SAP source) et \texttt{Extraction\_Timestamp} (horodatage de l'extraction), qui permettent de surveiller plusieurs instances SAP dans la même base et de reconstituer un historique.

Le schéma de chaque document est flexible (BSON sans contrainte de colonne), ce qui permet d'adapter la structure aux spécificités de chaque transaction sans migration de schéma. Les exemples suivants illustrent la structure réelle des enregistrements.

\begin{lstlisting}[caption={Enregistrement ST22 — erreur d'exécution ABAP}]
{
  "Current_Date"       : "2024-12-11",
  "Time"               : "14:46:45",
  "Application_Server" : "ECC-SYS-00",
  "User_Name"          : "USER____",
  "Client"             : "620",
  "Runtime_Error_Name" : "DBIF_RSQL_INVALID_REQUEST",
  "SID"                : "ECC",
  "Extraction_Timestamp": "2024-12-12 11:49:32"
}
\end{lstlisting}

\begin{lstlisting}[caption={Enregistrement SM37 — job annulé}]
{
  "Job"            : "Z3_WM_UPDATE",
  "Job_Created_By" : "BATCH_USER",
  "Status"         : "Cancelled",
  "Start_Date"     : "2025-04-24",
  "Start_Time"     : "19:08:01",
  "End_Date"       : "2025-04-24",
  "End_Time"       : "19:08:01",
  "SID"            : "ECC",
  "Extraction_Timestamp": "2025-04-25 10:59:14"
}
\end{lstlisting}

\subsection{Fichiers Produits par le Pipeline RAG}

L'extraction et la transformation des données documentaires produisent une famille de fichiers intermédiaires et finaux.

\begin{table}[H]
\centering
\caption{Fichiers produits par le pipeline de traitement documentaire}
\label{tab:fichiers}
\begin{tabularx}{\textwidth}{|p{5.8cm}|X|}
\hline
\textbf{Fichier} & \textbf{Contenu} \\
\hline
\texttt{output/<COURS>/structured\_content.json} & Contenu hiérarchique extrait par livre (sections, procédures, tableaux, métadonnées) \\
\hline
\texttt{normalized/unified\_content.json} & Documents normalisés issus de toutes les sources, après nettoyage et unification du schéma \\
\hline
\texttt{chunked\_for\_embedding.json} & Segments finaux après découpage hiérarchique (parents + enfants + singles) \\
\hline
\texttt{metadata.json} & Métadonnées des 76\,463 chunks filtrés, alignées par index avec les vecteurs \\
\hline
\texttt{embeddings.npy} & Vecteurs d'embeddings — forme (76\,463, 1\,024), float32 \\
\hline
\end{tabularx}
\end{table}

\subsection{Schéma de la Table PostgreSQL \texttt{sap\_chunks}}

Tous les segments documentaires sont stockés dans une table PostgreSQL dont le schéma supporte simultanément la recherche sémantique par vecteur, le filtrage par métadonnées et la recherche plein texte BM25.

\begin{lstlisting}[language=SQL, caption={Schéma de la table sap\_chunks}]
CREATE TABLE sap_chunks (
    chunk_id            UUID PRIMARY KEY,
    text                TEXT NOT NULL,
    title               TEXT,
    source_type         VARCHAR,     -- 'technical_book' | 'sap_help'
    source_book         TEXT,
    module              VARCHAR,     -- BASIS | ABAP | MM | SD | FI | SECURITY | HANA | FIORI
    section_type        VARCHAR,     -- 'section' | 'procedure' | 'note' | 'table'
    breadcrumb          TEXT[],      -- chemin hierarchique de la section
    mentioned_tcodes    TEXT[],      -- T-codes detectes (SM37, ST22, SE38...)
    mentioned_sap_notes TEXT[],      -- numeros de notes SAP (5 a 7 chiffres)
    is_parent           BOOLEAN,     -- TRUE si chunk parent (non indexe)
    parent_chunk_id     UUID,        -- reference au chunk parent
    text_search         TSVECTOR,    -- colonne pour la recherche full-text BM25
    embedding           vector(1024) -- vecteur BGE-large-en-v1.5
);

-- Index de recherche vectorielle approximative :
CREATE INDEX sap_chunks_hnsw ON sap_chunks
    USING hnsw (embedding vector_cosine_ops)
    WITH (m = 16, ef_construction = 64);

-- Index full-text BM25 :
CREATE INDEX sap_chunks_fts ON sap_chunks USING gin(text_search);

-- Index sur les T-codes pour le filtrage par transaction :
CREATE INDEX sap_chunks_tcodes ON sap_chunks USING gin(mentioned_tcodes);
\end{lstlisting}

L'index HNSW (\textit{Hierarchical Navigable Small World}) transforme la recherche par similarité cosinus d'une complexité linéaire $O(N)$ en une complexité approximative $O(\log N)$, ce qui rend le système viable à l'échelle de 76\,463 vecteurs de 1\,024 dimensions. Les paramètres \texttt{m=16} et \texttt{ef\_construction=64} équilibrent précision de rappel et temps de construction. L'index HNSW occupe 154\,Mo, l'index GIN sur les T-codes 11\,Mo.

% ─────────────────────────────────────────────────────────────────────────────
\section{Exploration des Données}
% ─────────────────────────────────────────────────────────────────────────────

\subsection{Corpus Documentaire RAG — Distribution Volumétrique}

Après extraction, normalisation, découpage et filtrage des chunks de figures, le corpus final compte \textbf{76\,463 segments} répartis sur deux types de sources.

\begin{table}[H]
\centering
\caption{Distribution des segments par type de source}
\label{tab:distribution-source}
\begin{tabular}{lrr}
\toprule
\textbf{Source} & \textbf{Segments} & \textbf{Part} \\
\midrule
Livres de formation (\texttt{technical\_book}) & 40\,899 & 53,5\,\% \\
Portail d'aide SAP (\texttt{sap\_help})        & 35\,564 & 46,5\,\% \\
\midrule
\textbf{Total} & \textbf{76\,463} & \textbf{100\,\%} \\
\bottomrule
\end{tabular}
\end{table}

\begin{table}[H]
\centering
\caption{Distribution des segments par module SAP}
\label{tab:distribution-module}
\begin{tabular}{lr}
\toprule
\textbf{Module SAP} & \textbf{Segments (approx.)} \\
\midrule
ABAP         & 12\,699 \\
Fiori        &  6\,263 \\
BASIS        &  4\,251 \\
HANA         &  3\,994 \\
MM / SD / FI & $\approx$ 8\,000 \\
Non classifié & $\approx$ 41\,256 \\
\bottomrule
\end{tabular}
\end{table}

\subsection{Structure du Découpage Hiérarchique}

Le pipeline de segmentation produit trois types de segments. Les \textit{singles} sont des documents courts tenant en un seul segment (73\,137 unités) ; ils sont directement indexés. Les \textit{parents} sont des sections longues conservées en version intégrale pour le contexte élargi (8\,795 unités) ; ils ne sont pas embedés mais servent de réserve lors de la génération. Les \textit{enfants} sont les fragments issus de la subdivision des sections longues (32\,194 unités) ; ils sont embedés et indexés pour la recherche. La taille moyenne d'un segment feuille (single ou enfant) est de 175 tokens.

\begin{table}[H]
\centering
\caption{Structure du corpus après segmentation hiérarchique}
\label{tab:chunking}
\begin{tabular}{lrr}
\toprule
\textbf{Type de segment} & \textbf{Quantité} & \textbf{Indexé dans pgvector} \\
\midrule
Singles (documents courts) & 73\,137 & Oui \\
Parents (sections longues) &  8\,795 & Non (contexte élargi) \\
Enfants (fragments)        & 32\,194 & Oui \\
\midrule
Total produit par le chunker & 114\,126 & --- \\
Total après filtrage figures & \textbf{76\,463} & \textbf{Oui} \\
\bottomrule
\end{tabular}
\end{table}

\subsection{Couverture des T-codes SAP}

Plus de 80 codes transaction SAP distincts sont détectés dans le corpus. Les cinq T-codes les plus fréquents sont SM37, ST22, SE38, SM21 et SM50, qui concentrent l'essentiel des occurrences. La distribution est très asymétrique : quelques transactions très documentées (monitoring Basis, développement ABAP) cumulent la majorité des mentions, tandis que la plupart des 80+ T-codes détectés n'apparaissent que dans quelques segments.

\subsection{Évaluation de la Récupération}

Une évaluation sur 18 requêtes de test couvrant quatre catégories — dépannage (5 requêtes), configuration (3), procédure (5), explication (5) — donne les métriques suivantes sur le corpus filtré :

\begin{table}[H]
\centering
\caption{Métriques de récupération sur le corpus de 76\,463 segments}
\label{tab:eval}
\begin{tabular}{llr}
\toprule
\textbf{Métrique} & \textbf{Signification} & \textbf{Score} \\
\midrule
Recall@5  & Part des segments pertinents dans les 5 premiers  & 0,5256 \\
Recall@10 & Part des segments pertinents dans les 10 premiers & 1,0000 \\
MRR       & Rang moyen du premier résultat pertinent           & 1,0000 \\
nDCG@5    & Qualité du classement dans les 5 premiers          & 0,8747 \\
nDCG@10   & Qualité du classement dans les 10 premiers         & 0,9460 \\
\bottomrule
\end{tabular}
\end{table}

La MRR de 1,0 indique qu'au moins un résultat pertinent apparaît systématiquement en première position pour chaque requête de test. Le Recall@10 de 1,0 confirme que tous les segments pertinents sont récupérés dans les dix premiers résultats. Le score nDCG@10 de 0,946 reflète une qualité de classement quasi-parfaite.

% ─────────────────────────────────────────────────────────────────────────────
\section{Vérification de la Qualité des Données}
% ─────────────────────────────────────────────────────────────────────────────

\subsection{Problèmes Détectés}

L'extraction automatique de texte depuis des PDFs techniques et le scraping de pages web exposent le pipeline à plusieurs catégories de bruit identifiées lors de l'analyse du corpus.

\paragraph{Artefacts CID (Character Identifier).}
Le format PDF stocke les glyphes de certaines polices via des identifiants internes appelés codes CID. Lorsque la police n'est pas embarquée dans le fichier, pdfplumber ne peut pas décoder ces glyphes et produit des séquences du type \texttt{(cid:123)}. Ces séquences n'ont aucune valeur sémantique et dégradent la qualité de l'embedding si elles ne sont pas supprimées.

\paragraph{Résidus HTML dans les chunks du portail SAP.}
Les pages extraites du portail SAP Help contiennent par endroits des balises HTML mal nettoyées (\texttt{<b>}, \texttt{<i>}, \texttt{\&nbsp;}, \texttt{\&amp;}). Ces résidus perturbent la tokenisation et introduisent du bruit dans les vecteurs d'embedding.

\paragraph{Faux positifs dans la détection de T-codes.}
L'heuristique de détection de T-codes par expression régulière (\texttt{[A-Z]\{2,4\}\textbackslash d\{1,3\}}) génère des faux positifs sur des abréviations comme \texttt{BC100} (numéro de cours SAP) ou des identifiants de section. Ces faux positifs créent des associations incorrectes entre un segment et un T-code inexistant.

\paragraph{Gap vocabulaire du modèle d'embedding.}
Le modèle \texttt{BAAI/bge-large-en-v1.5} est un modèle généraliste entraîné sur du texte académique et web. Il ne connaît pas nativement la signification des codes transaction SAP. Sans traitement, une requête contenant « SM37 » et un segment documentant « le moniteur de jobs en arrière-plan » ne seraient pas rapprochés par similarité cosinus, car les deux formulations sont lexicalement distantes dans l'espace vectoriel du modèle.

\paragraph{Chunks de figures.}
La version initiale du corpus comptait environ 26\,500 segments issus des légendes et descriptions de figures extraites des PDFs. Ces segments contiennent des informations très fragmentaires (légendes courtes, numéros de figure, textes alternatifs) qui dégradent la précision de la récupération en introduisant des résultats non pertinents.

\paragraph{Séquences de points issues des tables des matières.}
Les tables des matières des PDFs SAP génèrent des lignes du type \texttt{Chapitre 3 .......... 42} qui, une fois extraites, polluent les segments avec des suites de points sans valeur sémantique.

\paragraph{Doublons inter-sources.}
Certaines pages du portail d'aide SAP reproduisent des extraits de livres de formation. Sans déduplication, ces segments dupliqués introduisent un biais dans la récupération en sur-représentant certains contenus.

\subsection{Actions Appliquées}

\begin{table}[H]
\centering
\caption{Problèmes détectés et traitements appliqués}
\label{tab:qualite}
\begin{tabularx}{\textwidth}{|p{4.5cm}|X|}
\hline
\textbf{Problème} & \textbf{Traitement appliqué} \\
\hline
Artefacts CID & Filtrage par expression régulière \texttt{(cid:\textbackslash d+)} en post-extraction \\
\hline
Résidus HTML & Nettoyage strict par BeautifulSoup avec suppression des balises et décodage des entités HTML \\
\hline
Faux positifs T-codes & Application d'une liste blanche de 80+ T-codes SAP connus en post-traitement \\
\hline
Gap vocabulaire & Expansion de synonymes : description textuelle de chaque T-code concaténée au segment lors de l'indexation et à la requête lors de la recherche \\
\hline
Chunks de figures & Filtrage complet des segments dont \texttt{source\_type = 'figure'} avant l'ingestion en base \\
\hline
Séquences de points & Suppression par regex des séquences \texttt{\.{3,}} dans la phase de normalisation \\
\hline
Doublons inter-sources & Déduplication par empreinte MD5 des 300 premiers caractères, avec conservation de la copie aux métadonnées les plus riches \\
\hline
\end{tabularx}
\end{table}

\subsection{Actions Prévues}

Plusieurs améliorations restent à mettre en œuvre pour les prochaines itérations du corpus.

\begin{itemize}
    \item \textbf{Enrichissement du corpus SAP Help} : la couverture des notes de correction (SAP Notes) et des messages d'erreur ABAP sera augmentée par une collecte plus systématique des sections correspondantes du portail.
    \item \textbf{Rééquilibrage des modules} : la surreprésentation d'ABAP et Fiori par rapport aux modules FI et CO sera corrigée par l'ajout de livres dédiés lors de la prochaine campagne de collecte.
    \item \textbf{Validation humaine} : un échantillon de 500 segments représentatifs fera l'objet d'une relecture manuelle pour mesurer précisément le taux de bruit résiduel.
    \item \textbf{Modèle d'embedding domaine-spécifique} : le remplacement de \texttt{bge-large-en-v1.5} par un modèle fine-tuné sur la documentation SAP est en cours d'évaluation, en utilisant les paires requête/segment identifiées lors de la phase d'évaluation.
\end{itemize}
